\section{Ceros de funciones holomorfas}
\subsection{Ceros de funciones holomorfas}

\begin{definición} [Ceros de una función holomorfa]
    Sea $\Omega \subset \mathbb{C}$ un dominio, y sea $f : \Omega \to \mathbb{C}$ una función holomorfa en $\Omega$. Denotemos por $Z_f$ al conjunto de ceros de $f$ en $\Omega$, es decir:
    \[
        Z_f := \{ z \in \Omega : f(z) = 0 \}.
    \]
\end{definición}

\begin{definición} [Cero aislado]
    Diremos que $z_0 \in Z_f$ es un \textbf{cero aislado} de $f$ si existe $r > 0$ tal que:
    \[
        f(z) \neq 0 \quad \forall z \in D(z_0, r) \setminus \{ z_0 \}.
    \]
    o sea, si $z_0$ es un punto aislado de $Z_f$.
\end{definición}

\begin{proposición}
    Sea $\Omega \subset \mathbb{C}$ un dominio, y sea $f : \Omega \to \mathbb{C}$ una función holomorfa en $\Omega$. Entonces, si $z_0 \in Z_f$ es un cero aislado de $f$, existe $k \in \mathbb{N}$ y una función $g$ holomorfa en $\Omega$ tal que $g(z_0) \neq 0$ y \label{prop:ceros_aislados}
    \[
        f(z) = (z - z_0)^k g(z) \quad \forall z \in \Omega.
    \]
\end{proposición}

\begin{proof}
    Sea $z_0 \in \Omega$ y $f \in \mathcal{H}(\Omega)$, entonces por el teorema de Taylor 
    \[
        f(z) = \sum_{n=0}^{+\infty} c_n (z - z_0)^n \quad \forall z \in D(z_0, d(z_0, \partial \Omega))
    \]
    donde
    \[
        c_n = \frac{f^{(n)}(z_0)}{n!} \quad \forall n \in \mathbb{N} \cup \{0\}.
    \]
    Como $f(z_0) = 0$, se tiene que $c_0 = 0$, pero como $z_0$ es un punto aislado de $Z_f$, no se cumple que $c_n = 0$ para todo $n \in \mathbb{N}$ (pues en ese caso se daría que $f \equiv 0$ en $D(z_0, d(z_0, \partial \Omega))$). Así por la buena ordenación de los naturales, existe $k \in \mathbb{N}$ tal que:
    \[
        k = \min \{ n \in \mathbb{N} : c_n \neq 0 \}.
    \]
    Luego 
    \[
        f(z) = \sum_{n=k}^{+\infty} c_n (z - z_0)^n = (z - z_0)^k \sum_{n=0}^{+\infty} c_{n+k} (z - z_0)^n \quad \forall z \in D(z_0, d(z_0, \partial \Omega)).
    \]
    Si definimos 
    \[
        g(z) = \begin{cases}
            \frac{f(z)}{(z - z_0)^k} & \text{si } z \in \Omega \setminus \{ z_0 \}, \\
            c_k & \text{si } z = z_0,
        \end{cases}
    \]
    Tenemos que $g \in \mathcal{H}(\Omega \setminus \{ z_0 \})$ y $g$ es continua en $z_0$, luego $g \in \mathcal{H}(\Omega)$ por el teorema de prolongación analítica. Con lo cual
    \[
        f(z) = (z - z_0)^k g(z) \quad \forall z \in \Omega,
    \]
\end{proof}

\textbf{Nota:}
\begin{itemize}
    \item Al número $k$ del resultado anterior se le llama \textbf{orden del cero} $z_0$ de $f$, y corresponde a la multiplicidad del cero $z_0$ de $f$.
    \item Observemos que una función holomorfa $f$ en un dominio $\Omega$ con $f \neq 0$ no tiene ceros de orden infinito (caso que ocurre en $\mathbb{R}$).
    \item Si $z_0$ es un caso aislado de $f$, entonces $z_0$ es un cero aislado de orden $k$ si y sólo si $f^{(j)}(z_0) = 0$ para $j = 0, \ldots, k-1$.
\end{itemize}

\subsection{Principio de identidad}

\begin{teorema} [Principio de identidad]
    Sea $\Omega \subset \mathbb{C}$ un dominio, y sea $f \in \mathcal{H}(\Omega)$, si $Z_f$ tiene un punto de acumulación en $\Omega$, entonces $f \equiv 0$ en $\Omega$.    
\end{teorema}

\begin{proof}
    Sea 
    \[
        E := \{z \in \Omega : z \text{ es un punto de acumulación de } Z_f \}
    \]
    por hipótesis, $E \neq \emptyset$. Vamos a probar que $E$ y $\Omega \setminus E$ son abiertos en $\Omega$, con lo cual, como $\Omega$ es conexo, se tendrá que $\Omega \setminus E = \emptyset$, y por tanto $E = \Omega$.\\
    Pero todo elemento de $E$ es un cero de $f$, ya que si $z_0 \in E$, existe una sucesión $(z_n)_{n \in \mathbb{N}} \subset Z_f \setminus \{ z_0 \}$ tal que $z_n \to z_0$ cuando $n \to +\infty$, Asi
    \[
        \begin{cases}
            f \text{ es continua en } z_0, \\
            z_n \xrightarrow{n \to +\infty} z_0, 
        \end{cases} \implies f(z_0) = \lim_{n \to +\infty} f(z_n) = 0,
    \]
    pues $z_n \in Z_f$ para todo $n \in \mathbb{N}$. Luego $z_0 \in Z_f$.\\
    Con lo cual 
    \[
        E \subset Z_f \subset \Omega \subset E \implies Z_f = \Omega
    \]
    \[
        \implies f \equiv 0 \text{ en } \Omega.
    \]
    \begin{itemize}
        \item Sea $z_0 \in E \subset \Omega$, entonces por el teorema de Taylor, existe $(c_n)_{n = 0}^{\infty} \subset \mathbb{C}$ tal que:
        \[
            f(z) = \sum_{n=0}^{+\infty} c_n (z - z_0)^n \quad \forall z \in \underbrace{D(z_0, d(z_0, \partial \Omega))}_{D}.
        \]
        Como $z_0 \in E$, se tiene que $\exists (z_n)_{n \in \mathbb{N}} \subset Z_f \setminus \{ z_0 \}$ tal que $z_n \to z_0$ cuando $n \to +\infty$. Luego por el ejercicio 10 de la hoja 3 (Sea $f(z)=\sum_{n=0}^{+\infty} a_n z^n$ una función suma de una serie con radio de convergendia $R>0$; entonces si existe una sucesion $(z_n)_{n \in \mathbb{N}} \subset D(0,R)-\{0\}$ tal que $z_n \rightarrow 0$ y $f(a_n)=0 \forall n\in \mathbb{N}$ entonces $f=0 \quad \forall z \in D(0,R)$ ) deducimos que $f \equiv 0$ en $D$, luego $D \subset E$, con lo cual $E$ es abierto en $\Omega$.
        \item Si $z_0 \in \Omega \setminus E$, entonces
        \[
            \begin{cases}
                \text{o bien } z_0 \text{ es un cero aislado de } f, \\
                \text{o bien } f(z_0) \neq 0.
            \end{cases}
        \]
        en cualquier caso existe $r > 0$ tal que en $D(z_0, r) \setminus \{ z_0 \}$ $f$ no se anula, por tanto $D(z_0, r) \subset \Omega \setminus E$, con lo cual $\Omega \setminus E$ es abierto en $\Omega$.
    \end{itemize}
\end{proof}

\begin{teorema} [Principio de identidad - reformulación]
    Sea $\Omega \subset \mathbb{C}$ un dominio, y sean $f, g \in \mathcal{H}(\Omega)$. Si el conjunto
    \[
        S := \{ z \in \Omega : f(z) = g(z) \}
    \]
    tiene un punto de acumulación en $\Omega$, entonces $f \equiv g$ en $\Omega$.
\end{teorema}


\ejemplo{

    \begin{enumerate}
        \item Sea $f(1/n) = \sin(\frac{1}{n})$ para $n \in \mathbb{N}$ siendo $f \in \mathcal{H}(\mathbb{C})$. Entonces nos preguntamos si tiene solución única. La respuesta es afirmativa siendo $f(z) = \sin(z) \ \forall z \in \mathbb{C}$ pues, $\{\frac{1}{n} : n \in \mathbb{N}\}$ tiene un único punto de acumulación, el $0 \in \mathbb{C}$.
        \item Lo mismo nos preguntamos para la función $f(\frac{1}{n \pi}) = 0$ para $n \in \mathbb{N}$ siendo $f \in \mathcal{H}(\mathbb{C} \setminus \{0\})$. En este caso, la respuesta es negativa, ya que $f \equiv 0$ en $\mathbb{C} \setminus \{0\}$ y $g(z) = \sin(\frac{1}{z})$ en $\mathbb{C} \setminus \{0\}$ son dos soluciones. El conjunto $\{\frac{1}{n \pi} : n \in \mathbb{N}\}$ tiene un único punto de acumulación, el $0 \notin \mathbb{C} \setminus \{0\}$.
    \end{enumerate}
}

\begin{observación}
    \vspace{-2em}
    \begin{itemize}
        \item El resultado anterior nos dice que una función holomorfa en un dominio $\Omega$ de $\mathbb{C}$ está unívocamente determinada por sus valores en un subconjunto $S \subset \Omega$ que tenga un punto de acumulación en $\Omega$. Esto no ocurre en el caso real.
        \item Las funciones $e^{z}$, $\sin(z)$ y $\cos(z)$ son las úncias extensiones holomorfas de $e^{x}$, $\sin(x)$ y $\cos(x)$ respectivamente en un abierto conexo de $\mathbb{C}$ que contenga a $\mathbb{R}$.
        \item Si $f \in \mathcal{H}(\Omega)$ siendo $\Omega \subset \mathbb{C}$ un dominio y si $f \not\equiv  0$, entonces el conjunto $Z_f$ de ceros de $f$ en $\Omega$ no tiene puntos de acumulación en $\Omega$, luego todos los puntos de $Z_f$ son ceros aislados de $f$ y, por tanto, $Z_f$ es un conjunto finito o numerable.\\
        Esto último se debe a que en cada $K \subset \Omega$ compacto, sólo puede haber un número finito de ceros de $f$, ya que
        \[
            \Omega = \bigcup_{n=1}^{\infty} K_n \quad \text{ con } K_n \subset \Omega \text{ compacto } \forall n \in \mathbb{N}
        \]
        basta considerar para cada $n \in \mathbb{N}$
        \[
            K_n := \overline{D}(0,n) \cap \left\{ z \in \Omega : d(z, \partial \Omega) \geq \frac{1}{n} \right\}
        \]
        \item Se puede probar fácilmente que 
        \[
            \sin^2(z) + \cos^2(z) = 1 \quad \forall z \in \mathbb{C}
        \]
        pues se cumple en $\mathbb{R}$, que tiene puntos de acumulación en $\mathbb{C}$, luego se cumple en todo $\mathbb{C}$.
    \end{itemize}
\end{observación}

\subsection{Principio del módulo máximo}

\begin{teorema} [Principio del módulo máximo]
    Sea $\Omega \subset \mathbb{C}$ un dominio, y sea $f \in \mathcal{H}(\Omega)$. Si $|f|$ tiene un máximo relativo en un punto $z_0 \in \Omega$, entonces $f$ es constante en $\Omega$.
\end{teorema}

\begin{proof}
    Como $z_0 \in \Omega$ es un punto de máximo relativo de $|f|$, $\Omega$ es abierto, existe $r_0 > 0$ tal que $D(z_0, r_0) \subset \Omega$ y:
    \[
        |f(z)| \leq |f(z_0)| \quad \forall z \in D(z_0, r_0).
    \]
    Para cada $r \in (0, r_0)$, aplicando la fórmula integral de Cauchy en la circunferencia $C(z_0, r)$, tenemos que:
    \[
        \left|f(z_0)\right| = \left| \frac{1}{2 \pi i} \int_{C(z_0, r)} \frac{f(z)}{z - z_0} \, dz \right| = \frac{1}{2 \pi i} \left| \int_{0}^{2 \pi} \frac{f(z_0 + r e^{i t})}{r e^{i t}} i r e^{i t} \, dt \right| = \frac{1}{2 \pi} \left| \int_{0}^{2 \pi} f(z_0 + r e^{i t}) \, dt \right| 
    \]
    \[
        \leq \frac{1}{2 \pi} \int_{0}^{2 \pi} |f(z_0 + r e^{i t})| \, dt \leq \frac{1}{2 \pi} \int_{0}^{2 \pi} |f(z_0)| \, dt = |f(z_0)|.
    \]
    Usando como parametrización de la curva $\gamma(t) = z_0 + r e^{i t}$ para $t \in [0, 2 \pi]$.\\
    Por tanto,
    \[
        \frac{1}{2 \pi} \int_{0}^{2 \pi} |f(z_0 + r e^{i t})| \, dt = \frac{1}{2 \pi} \int_{0}^{2 \pi} |f(z_0)| \, dt
    \]
    con lo cual
    \[
        \int_{0}^{2 \pi} \underbrace{ \left( |f(z_0)| - |f(z_0 + r e^{i t})| \right)}_{\varphi(t)} dt = 0.
    \]
    como $\varphi(t) = |f(z_0)| - |f(z_0 + r e^{i t})|$ es continua positiva y $\int_{0}^{2 \pi} \varphi(t) dt = 0$, $\varphi(t) \equiv 0$ en $[0, 2 \pi]$, es decir,
    \[
        |f(z_0 + r e^{i t})| = |f(z_0)| \quad \forall t \in [0, 2 \pi]
        \iff |f(z)| = |f(z_0)| \quad \forall z \in C(z_0, r).
    \]
    Como $r \in (0, r_0)$ es arbitrario, se tiene que $|f(z)| = |f(z_0)|$ para todo $z$ en el disco abierto $D(z_0, r_0)$.\\
    Luego $|f|$ es constante en $D(z_0, r_0)$, y también lo es en $\Omega$ por el principio de identidad, luego $f$ es constante en $\Omega$ por (2) de la \cref{prop:constantes_holomorfas}.
\end{proof}

\begin{corolario}
    Sea $\Omega \subset \mathbb{C}$ un dominio acotado y sea $f \in \mathcal{H}(\Omega)$ continua en $\overline{\Omega}$, entonces $|f|$ alcanza su máximo en $\partial \Omega$.
\end{corolario}

\begin{proof}
    Si $f$ es contante en $\Omega$, entonces $|f|$ es constante en $\overline{\Omega}$ por la continuidad de $f$, y alcanza su máximo en cualquier punto de $\partial \Omega$.\\
    Si $f$ no es constante en $\Omega$, entonces por el principio del módulo máximo, $|f|$ no puede alcanzar su máximo en ningún punto de $\Omega$, luego debe alcanzarlo en $\partial \Omega$, pues $\overline{\Omega}$ es compacto al ser cerrado y actodo (luego debe alcanzar sus valores máximo y mínimo).
\end{proof}

Podemos usar el principio del módulo máximo para demostrar el teorema Schwarz.

\subsection{Lema de Schwarz}

\begin{lema} [Lema de Schwarz]
    Sea $D := D(0,1)$ y sea $f \in \mathcal{H}(D)$ tal que:
    \begin{itemize}
        \item $|f(z)| \leq 1$ para todo $z \in D$,
        \item $f(0) = 0$.
    \end{itemize}
    Entonces:
    \begin{enumerate}
        \item $|f(z)| \leq |z|$ para todo $z \in D$,
        \item $|f'(0)| \leq 1$,
    \end{enumerate}
    Además, si para algún $z_0 \in D \setminus \{0\}$ se cumple que $|f(z_0)| = |z_0|$ o si $|f'(0)| = 1$, entonces $f$ es una rotación, por lo que existe $\theta \in \mathbb{R}$ tal que:
    \[
        f(z) = e^{i \theta} z \quad \forall z \in D.
    \]
\end{lema}

\begin{proof}
    Consideremos la función:
    \[
    g(z) = 
    \begin{cases}
        \frac{f(z)}{z} & \text{si } z \in D \setminus \{0\}, \\
        f'(0) & \text{si } z = 0
    \end{cases}
    \implies g \in \mathcal{H}(D).
    \]
    Por lo que $g$ es continua en $0$ pues:
    \[
    \lim_{z \to 0} g(z) =
    \lim_{z \to 0} \frac{f(z)}{z} =
    \lim_{z \to 0} \frac{f(z) - f(0)}{z - 0} = 
    f'(0) =
    g(0)
    \]
    Luego por el principio de prolongación, $g \in \mathcal{H}(D)$.\\
    Si $z_1 \in D \setminus \{0\}$, entonces para cada $r \in (|z_1|, 1)$, se tiene que por el principio del módulo máximo:
    \[
    |g(z_1)| \leq 
    \sup_{|z|=r} |g(z)| =
    \sup_{|z|=r} \frac{|f(z)|}{|z|} =
    \sup_{|z|=r} \frac{|f(z)|}{r} \leq
    \frac{1}{r}
    \]
    Así, $g(z_1) \leq 1 \, \quad \forall z_1 \in (z_1|, 1)$, y por tanto:
    \[
    |g(z_1)| \leq \lim_{r \to 1^-} \frac{1}{r} = 1
    \]
    Por lo que $|g(z)| \leq 1$ para todo $z \in D \setminus \{0\}$, y por continuidad, $|g(0)| \leq 1$, con lo cual $\boxed{|g(z)| \leq 1 \quad \forall z \in D}$, que equivale a que $|f(z)| \leq |z| \quad \forall z \in D$ y $|f'(0)| \leq 1$.\\
    Si $\exists z_0 \in D \setminus \{0\}$ tal que $|f(z_0)| = |z_0|$, entonces $|g(z_0)| = 1$, luego $|g|$ tiene un maximo local en $z_0 \in D$, y por el principio del módulo máximo, $g$ es constante en $D$, luego $\exists c \in \mathbb{C}$ con $|c| = 1$ tal que $g(z) = c$ para todo $z \in D$, con lo cual $f(z) = c z$ para todo $z \in D$.\\
    Ahora, sea $\theta_0 \in \mathbb{R}$ tal que $e^{i \theta_0} = c$, entonces $g(z) = e^{i \theta_0} $ para todo $z \in D$, con lo cual $f(z) = e^{i \theta_0} z$ para todo $z \in D$.\\
    Finalmente, si $|f'(0)| = 1$, entonces $|g(0)| = 1$ y repetimos el resultado anterior.
\end{proof}

El lema de Schwarz tiene varias situaciones mas generales, como el siguiente corolario.

\begin{corolario}
    Si $f \in \mathcal{H}(D)$ y $f(D) \subset \overline{D}(0,M)$ para algún $M > 0$, aplicamos el lema de Schwarz a la función $\frac{f}{M}$, obteniendo que $|f(z)| \leq M |z| \quad \forall z \in D$ y que $|f'(0)| \leq M$.
\end{corolario}

\begin{corolario}
    Si tenemos $f \in \mathcal{H}(D(0,R) = D_R)$ y $f(D_R) \subset \overline{D}(0,M)$ para algún $M, R > 0$, aplicamos el corolario anterior y obtenemos $g(z) = f(R z) \quad \forall z \in D$, luego:
    \[
    |f(z)| = 
    |f(R \frac{z}{R})| =
    |g(\frac{z}{R})| \leq
    M |\frac{z}{R}| =
    \frac{M}{R} |z| \quad \forall z \in D
    \]
\end{corolario}
